\begin{abstract}
Host-to-device MMIO direct push---issuing posted 64-byte writes such as \texttt{MOVDIR64B} directly to NIC or accelerator memory---promises to bypass host last-level cache (LLC) and DRAM traffic while delivering high message rates from a single core. On a dual-target testbed combining an Intel Sapphire Rapids Xeon host with both an off-chip PCIe Gen5\,$\times$4 Marvell OCTEON~10 CN106 SmartNIC and an on-die Intel DSA accelerator, we show that this promise holds only under strict pacing: a single core pushing 175\,Mpps to SmartNIC \texttt{BAR4} degrades co-located host workloads by under 1\%, whereas traditional host-memory DMA rings cut co-located 6-core STREAM bandwidth by 9.9\%--41.0\%.

Without pacing, however, posted MMIO direct push triggers two severe failure modes that stem from the same open-loop backpressure chain. First, across both off-chip PCIe \texttt{BAR4} (32,768 cache lines) and an on-die 4\,KiB (64-line) Intel DSA portal, unpaced multi-core writes saturate the egress \texttt{IRP}/\texttt{M2PCIe} queues and spill into the socket's shared 32-slice \texttt{CHA} Table of Requests (\texttt{TOR}). Crossing a universal 330--430-entry \texttt{WCiLF} cliff collapses co-located 6-core STREAM bandwidth by 67.0\%--96.7\% and inflates 16\,MiB L3 tail latency by 7.7$\times$--116$\times$---even on the 4\,KiB DSA portal at zero DMA and zero drops, and even when the accelerator discards 94.6\% of descriptors in hardware. Second, on the device side, unpaced producer--consumer frontier collisions collapse throughput by 5$\times$--10$\times$ at two distinct levels: under default \texttt{Normal-WB} pages, active receiver polling locks the inbound bridge at 13.1--13.5\,Mmsg/s via L1/L2 snoop invalidation; remapping the ring with standard Linux \texttt{MT\_NORMAL\_NC} (\texttt{0x44}) eliminates private-cache snoops and DRAM reads to reach 160.1--174.1\,Mpps, yet over-provisioned polling still triggers L3 \texttt{EXLMN} collisions (16.3--22.8\,Mmsg/s) and overflows the device's $\sim$387-slot posted buffer right back into the host \texttt{CHA TOR} cliff (428--521 entries).

We distill pacing and receiver-consumption principles that close both halves of this loop. Coupling aggregate rate pacing with micro-burst step bounding (\texttt{b\,=\,128}) restores 0\% neighbor degradation at 8-core 165\,Mpps on \texttt{BAR4}, while a hardware-enforced disjoint-line rule ($\sum_i K_i < 200$) bounds \texttt{TOR} occupancy without software flow control. On a dual-host 100\,GbE testbed, a paced direct-push datapath increases single-core saturation throughput by 55\% on an in-memory KV store and 22.5\% on unmodified upstream Memcached~1.6.45 (Facebook ETC), cuts small-request bidirectional PCIe traffic by 59.4\%--62.5\%, and reduces net active host CPU cycles by 27.2\%--32.2\%.
\end{abstract}
